\begin{lemma}[Existence of a maximum-defective choice]
Let \(G\) be a finite loopless multigraph, let \(e_0\) be an edge, let
\(T\) be a tree-sequence with first edge \(e_0\), let
\(C\subseteq [k]\), and let \(\varphi\) be a proper partial
\(k\)-edge-coloring leaving exactly \(e_0\) uncolored.  If \(V(T)\) is
closed but not strongly closed under \(\varphi\), then there exists an
admissible maximum-defective choice for
\((G,e_0,T,C,\varphi)\) in the sense of
\noderef{MaximumDefectiveChoice}.  In particular, the choice retains an
exact defective boundary edge identity, its displayed incident endpoint,
its other endpoint outside \(V(T)\), the symmetric endpoint equation, and
its exact color.
\end{lemma}
\begin{proof}
Put \(X=V(T)\).  Since \(X\) is closed but is not strongly closed,
the second clause of \noderef{ColoredVertexSetStronglyClosed} fails.
Consequently there are distinct exact boundary edge identities \(f,g\)
and a color \(c\) such that
\[
 f,g\in\partial_G(X),\qquad \varphi(f)=c=\varphi(g).             \tag{1}
\]
Both edges therefore belong to
\(\partial_{\varphi,c}(X)\), as defined by
\noderef{BoundaryColorEdges}.  The two-element set \(\{f,g\}\) is a
subset of that boundary-color set, so its cardinality is at least two.
Thus \(c\) is a defective color by \noderef{DefectiveColor}, and each of
\(f,g\) is a defective edge by \noderef{DefectiveEdge}.

Unpack the membership \(f\in\partial_G(X)\) using
\noderef{BoundaryEdges}.  It supplies vertices \(v\in X\) and
\(u\notin X\) with
\[
                         \operatorname{ends}_G(f)=\{v,u\}.       \tag{2}
\]
In particular \(f\) is incident with \(v\) in the sense of
\noderef{Incident}.  Hence \(v\) is a defective vertex under
\(\varphi\) by \noderef{DefectiveVertex}.

We next perform the finite maximum in the manuscript footnote.  Let
\(I\) be the set of displayed positions \(a\) of \(T\) for which there
exist a partial coloring \(\pi\) satisfying all four conditions
\[
 \pi\text{ is proper},\quad \pi\text{ leaves exactly }e_0
 \text{ uncolored},\quad
 \pi\text{ is }(T,C,\varphi)\text{-stable},\quad
 T(a)\text{ is defective under }\pi.                            \tag{3}
\]
The displayed position set is finite by \noderef{TreeSequence}.  It is
nonempty: the vertex \(v\) obtained above has a unique displayed position
because \(v\in X=\noderef{TreeSequenceVertexSet}(T)\), and the identity
coloring \(\pi=\varphi\) satisfies (3).  Its stability is immediate from
both clauses of \noderef{StableColoring}.

Choose the greatest position \(a_*\in I\), and choose a coloring
\(\pi_*\) witnessing its membership in \(I\).  Set
\(v_*=T(a_*)\).  If \(\rho\) is any other proper coloring leaving
\(e_0\) uncolored, stable for \((T,C,\varphi)\), and \(w\) is a
defective vertex under \(\rho\), then
\(w=T(b)\) for a unique displayed position \(b\).  The four defining
conditions put \(b\) in \(I\), so \(b\leq a_*\).  If equality holds,
injectivity of the displayed vertex map gives \(w=v_*\); otherwise
\(b<a_*\), which is precisely
\(w\prec_T v_*\) by \noderef{TreeVertexOrder}.  This proves the required
global maximality over both colorings and their defective vertices.

Finally unpack \noderef{DefectiveVertex} for \(v_*\) under \(\pi_*\).
Choose the resulting exact defective edge \(e\) incident with \(v_*\).
Unpacking \noderef{DefectiveEdge} gives its exact color
\(\delta\) and \(e\in\partial_{\pi_*,\delta}(X)\), hence
\(e\in\partial_G(X)\) and \(\pi_*(e)=\delta\).  Unpacking the latter
boundary membership gives an inside endpoint \(x\in X\), an outside
endpoint \(y\notin X\), and
\(\operatorname{ends}_G(e)=\{x,y\}\).  Since \(e\) is incident with
\(v_*\in X\), equation of the unordered endpoint pair says that
\(v_*=x\) or \(v_*=y\); the second alternative contradicts
\(y\notin X\).  Thus \(v_*=x\), and
\[
 \operatorname{ends}_G(e)=\{v_*,y\},\qquad
 y\notin X,\qquad \pi_*(e)=\delta .
\]
Together with the properties in (3) and the proved maximality, these are
exactly all fields of \noderef{MaximumDefectiveChoice}.  They also retain
the identity of \(e\); no edge is replaced merely because another edge
has the same endpoints.
\end{proof}
